\documentclass[10pt,a4paper]{amsart}
\usepackage[margin=19mm]{geometry}
\usepackage{amsmath,amssymb,mathtools}
\usepackage[scr=rsfs,cal=euler]{mathalfa}
\usepackage{xfrac}
\usepackage[unicode,hidelinks,bookmarks=false]{hyperref}
\newcommand{\ie}{i.e.\ }
\newcommand{\cf}{cf.\ }
\newcommand{\resp}{respectively\ }
\newcommand{\Subsection}{\subsection}
\providecommand{\nobreakdash}{\nobreak\hbox{-}}
\setlength{\emergencystretch}{2em}
\multlinegap0pt
\usepackage{bm}
\usepackage{bbm}
\usepackage{dsfont}
\usepackage{xspace}
\usepackage{cancel}
\usepackage{mathtools}
\usepackage{amsthm}
\makeatletter
\@ifundefined{lemma}{\newtheorem{lemma}{Lemma}}{}
\makeatother
\newcommand{\R}{{\mathbb R}}
\newcommand{\const}{{\rm const}}
\newcommand{\supp}{\mathop{\rm supp}}
\title[On linear evolutionary equations with skew symmetric spatial]{On linear evolutionary equations with skew symmetric spatial operators}
\author{Evgeny Yu. Panov}
\date{Source: arXiv:2507.09314v2 (CC BY 4.0)}
\begin{document}
\maketitle

\noindent\emph{Source and licence.} On linear evolutionary equations with skew symmetric spatial operators. Version arXiv:2507.09314v2 (CC BY 4.0), distributed under the Creative Commons Attribution 4.0 International licence, as recorded in the arXiv licence metadata for this exact version and shown on the abstract page https://arxiv.org/abs/2507.09314v2. Full rights notice: licenses/arxiv-2507.09314-CC-BY-4.0.txt.\par\smallskip

\noindent\textit{Editorial scope.} Counted unit: the Lemma whose source label is \texttt{lmDL} in the article (source0.tex lines 567--574, source label \texttt{lmDL}), delivered together with the article's own complete original proof of it (source0.tex lines 576--645, one \texttt{proof} environment of 70 lines). No mathematics is written, repaired or re-derived: statement and proof are the article's own text line for line. Required context retained uncounted: lip (lines 305--307). Every definition, display and label of those lines is retained unchanged. Omitted from this package and not redistributed: 1--304, 308--566, 575--575, 646--722. Nothing inside a retained excerpt is omitted or altered: every retained body is delivered exactly as the article's lines are written, so no omission receipt of any kind accompanies this package. Citations: the retained excerpts carry no citation command at all, so no bibliography entry of the article is transcribed and no reference is redistributed. Reference adaptation: none was needed and none was made; every reference inside the delivered unit resolves inside it, so no unresolved reference or citation is produced. Printed slips kept literal: no printed word, symbol, hypothesis or number of the counted statement or of its proof was altered. Layout and class: the article's own class is \texttt{article} with packages including amsmath, amssymb, amsthm; the wrapper is the lane's pinned \texttt{amsart} editorial wrapper at 10pt on A4, which restates the theorem-like environments the retained text needs and adds the article's own macro definitions that the retained text needs (\texttt{\textbackslash R}, \texttt{\textbackslash const}, \texttt{\textbackslash supp}), with extra wrapper packages (bm, bbm, dsfont, xspace, cancel, mathtools). The delivered numbering is the unit's own. Assets: no retained span contains a figure environment, a graphics-inclusion command, a PDF-page inclusion, a table or a TikZ picture; no image file is copied into this package, the package declares no asset dependency and no image file is redistributed with it. Licence: version arXiv:2507.09314v2 of this article is distributed under the Creative Commons Attribution 4.0 International licence, as recorded in the arXiv licence metadata for this exact version and shown on the abstract page https://arxiv.org/abs/2507.09314v2. A full rights notice is delivered at licenses/arxiv-2507.09314-CC-BY-4.0.txt. Characters: every retained excerpt is pure ASCII, so the delivered engine, XeLaTeX with its default Latin Modern text font, reports no missing character.
\begin{equation}\label{lip}
|a(x)-a(y)|\le m|x-y| \ \forall x,y\in\R^n.
\end{equation}

\begin{lemma}\label{lmDL}
Let $a(x)\in C(\R^n)$ satisfy the global Lipschitz condition (\ref{lip}), that is
$|a(x)-a(y)|\le m|x-y|$ for all $x,y\in\R^n$; $u(x)\in L^p(\R^n)$, $1\le p<\infty$. Then
\begin{equation}\label{DL}
\frac{\partial}{\partial x_j}(au_k-(au)_k)(x)\to 0 \ \mbox{ in } L^p(\R^n)
\end{equation}
for all $j=1,\ldots,n$.
\end{lemma}

\begin{proof}
Since \[(au_k-(au)_k)(x)=k^n\int_{\R^n}(a(x)-a(y))\rho(k(x-y)) u(y)dy,\]
there exists the generalized derivatives
\begin{align}\label{rel}
\frac{\partial}{\partial x_j}(au_k-(au)_k)(x)=k^n\int_{\R^n}a_{x_j}(x)\rho(k(x-y)) u(y)dy+\nonumber\\ k^{n+1}\int_{\R^n}(a(x)-a(y))\rho_{z_j}(k(x-y)) u(y)dy=I_{1k}(x)+I_{2k}(x).
\end{align}
The first term in this sum
\begin{equation}\label{rel1}
I_{1k}(x)=a_{x_j}(x)\int_{\R^n}\rho_k(x-y) u(y)dy=a_{x_j}(x)u_k(x)\to a_{x_j}(x)u(x)
\end{equation}
as $k\to\infty$ in $L^p(\R^n)$ because $u_k\mathop{\to}\limits_{k\to\infty} u$ in $L^p(\R^n)$ by the known property of averaged functions while the derivative $a_{x_j}(x)\in L^\infty(\R^n)$ in view of the Lipschitz condition.
Next, we estimate the term $I_{2k}(x)$. Obviously,
\begin{align}\label{es1}
|I_{2k}(x)|\le k^{n+1}\int_{\R^n}|a(x)-a(y)||\rho_{z_j}(k(x-y))| |u(y)|dy\le \nonumber\\ \omega_k(x)\doteq mk^n\int_{\R^n}k|x-y||\rho_{z_j}(k(x-y))| |u(y)|dy.
\end{align}
By the properties of averaged functions $\omega_k(x)\in L^p(\R^n)$ and
\begin{equation}\label{rel2}
\omega_k(x)\to C|u(x)|
\end{equation}
as $k\to\infty$ both in $L^p(\R^n)$ and a.e. in $\R^n$. Here $\displaystyle C=m\int_{\R^n} |z||\rho_{z_j}(z)|dz=\const$.
Now, let $x$ be a common Lebesgue point of the vector $\nabla a(y)$ and the function $u(y)$. Then
\begin{align}\label{rel2-1}
I_{2k}(x)=k^{n+1}\int_{\R^n}(a(x)-a(y))\rho_{z_j}(k(x-y))(u(y)-u(x))dy+\nonumber\\ u(x)k^{n+1}\int_{\R^n}(a(x)-a(y))\rho_{z_j}(k(x-y))dy.
\end{align}
The first term in the right-hand part of (\ref{rel2-1}) is estimated as
\begin{align}\label{est2-1}
k^{n+1}\left|\int_{\R^n}(a(x)-a(y))\rho_{z_j}(k(x-y))(u(y)-u(x))dy\right|\le \nonumber\\
mk^n\int_{\R^n}k|x-y||\rho_{z_j}(k(x-y))||u(y)-u(x)|dy\mathop{\to}_{k\to\infty} 0
\end{align}
since $x$ is a Lebesgue point of $u(y)$. We introduce the function ${\displaystyle J_k(x)=k^{n+1}\int_{\R^n}(a(x)-a(y))\rho_{z_j}(k(x-y))dy}$ and represent it in the form
\begin{align}\label{rel2-2}
J_k(x)=k^{n+1}\int_{\R^n}(a(x)-a(y)-\nabla a(x)\cdot(x-y))\rho_{z_j}(k(x-y))dy+\nonumber\\ k^{n+1}\int_{\R^n}\nabla a(x)\cdot (x-y)\rho_{z_j}(k(x-y))dy=\nonumber\\
k^{n+1}\int_{\R^n}\int_0^1 (\nabla a(x+s(y-x))-\nabla a(x))\cdot(x-y)\rho_{z_j}(k(x-y))dsdy+\nonumber\\ \nabla a(x)\cdot\int_{\R^n}z\rho_{z_j}(z)dz.
\end{align}
We utilize here the identity
\[
a(x)-a(y)=\int_0^1\nabla a(x+s(y-x))\cdot(x-y)ds,
\]
which holds for a.e. $y\in\R^n$.
To estimate the first term in the right-hand part of (\ref{rel2-2}), we make the change of variables $z=s(x-y)$, resulting in
\begin{align}\label{est2-2}
k^{n+1}\left|\int_{\R^n}\int_0^1 (\nabla a(x+s(y-x))-\nabla a(x))\cdot(x-y)\rho_{z_j}(k(x-y))dsdy\right|\le\nonumber\\
k^n\int_{\R^n} |\nabla a(x-z)-\nabla a(x)|\rho_1(kz)dz,
\end{align}
where we denote
\[
\rho_1(y)=|y|\int_0^1s^{-n-1}|\rho_{z_j}(y/s)|ds.
\]
Remark that $\rho_1\ge 0$, $\supp\rho_1\subset B_1(0)$, $\displaystyle\int_{\R^n}\rho_1(y)dy=\int_{\R^n}|z||\rho_{z_j}(z)|dz$, and since $x$ is a Lebesgue point of the vector $\nabla a(y)$, we find
\[k^n\int_{\R^n} |\nabla a(x-z)-\nabla a(x)|\rho_1(kz)dz\mathop{\to}_{k\to\infty} 0\].
It follows from (\ref{rel2-2}) and (\ref{est2-2}) that
\begin{equation} \label{rel3}
J_k(x)\mathop{\to}_{k\to\infty} \nabla a(x)\cdot\int_{\R^n}z\rho_{z_j}(z)dz=-a_{x_j}(x).
\end{equation}
Here, we apply the integration by parts formula
\[\int_{\R^n}z\rho_{z_j}(z)dz=-\int_{\R^n}\frac{\partial z}{\partial z_j}\rho(z)dz=-e_j, \]
where $e_j$ is the $j$th basis vector in $\R^n$.
In view of (\ref{rel2-1}), (\ref{est2-1}) it follows from (\ref{rel3}) that $I_{2k}(x)\to -a_{x_j}(x)u(x)$ as $k\to\infty$ a.e. in $\R^n$. Further, we notice that
\[
|I_{2k}(x)+a_{x_j}(x)u(x)|^p\le 2^{p-1}(|I_{2k}(x)|^p+|a_{x_j}(x)u(x)|^p)\le 2^{p-1}((\omega_k(x))^p+|a_{x_j}(x)u(x)|^p).\]
The left-hand side of this inequality converges to zero a.e. in $\R^n$, while its right-hand side converges both in  $L^1(\R^n)$ and a.e. in $\R^n$, by relation (\ref{rel2}). Applying Fatou's lemma to the sequences
\[
2^{p-1}((\omega_k(x))^p+|a_{x_j}(x)u(x)|^p)\pm |I_{2k}(x)+a_{x_j}(x)u(x)|^p,
\]
we derive that
\[
\lim_{k\to\infty} \int_{\R^n}|I_{2k}(x)+a_{x_j}(x)u(x)|^pdx=0,
\]
that is, $\displaystyle I_{2k}(x)\mathop{\to}_{k\to\infty} -a_{x_j}(x)u(x)$ in $L^p(\R^n)$. This, together with  (\ref{rel}), (\ref{rel1}), gives the required relation (\ref{DL}). The proof is complete.
\end{proof}

\end{document}
